% Resumo do relatório -- EDITÁVEL À MÃO (exportado em 2026-11-05; specs/2026-11-05_versao_final_site R5).
% Resumido em 2026-11-05 para caber numa página (pedido do usuário); mesmo conteúdo do resumo do site, mais enxuto.
\setlength{\absparsep}{12pt}
\begin{resumo}
O Diagnóstico da Primeira Infância, desenvolvido pela Coordenadoria de Pesquisa, Avaliação e Políticas Públicas do Instituto Pereira Passos em parceria com a Casa Civil, apoia o desenho, a implementação e o monitoramento da nova Política Integrada da Primeira Infância. É um produto vivo, que deve ser atualizado à medida que surgirem novos dados ou novas necessidades; esta versão em relatório reproduz o conteúdo do site oficial, com mudanças apenas de organização e visual.

A principal tendência é a clara redução da participação das crianças até 72 meses (0 a 6 anos) na população, ligada sobretudo à queda dos nascimentos, que deve continuar a afetar os indicadores nos próximos anos. Os dados também mostram melhorias na prevenção de mortes evitáveis, a manutenção de desigualdades significativas entre as regiões da cidade e a persistência de uma parcela alta da população-alvo em territórios com muita violência ou em privação socioeconômica.

Produzir estatísticas e análises abaixo do nível municipal é sempre um desafio: a maior parte da produção estatística brasileira tem o município como menor unidade, e as exceções, como os Censos Demográficos, são esporádicas. Por isso, os dados vêm com frequência de fontes distintas, não concebidas para comparação, ou de registros administrativos que não se destinavam à análise.

Por fim, os processos sociais apontados podem ser tanto causa quanto consequência de outras variáveis, econômicas, tecnológicas ou demográficas, e o período inclui as crises da Zika (2015-2016) e da Covid-19 (2020-2021) e a recessão de 2015-2016. Assim, recomendamos que os dados sirvam como um mapa para orientar indicadores sensíveis a cada contexto, e não como metas em si mesmas.

\noindent\textbf{Palavras-chave}: Primeira infância. Indicadores sociais. Políticas públicas. Rio de Janeiro (RJ).
\end{resumo}
